%==============================================================================
%  An Exact Graded Curvature Tower for the Fractional Laplacian
%  Operator Realization and Wiener Chaos  (positive range 0 < H < 1)
%
%  Standalone sequel to Part I (fractional-Laplacian Bakry--Emery curvature).
%  No iterated-Brownian / pseudo-process content: that is the H>1 sequel.
%
%  Build:  pdflatex iterated_gamma_curvature_tower.tex   (twice for refs)
%==============================================================================
\documentclass[11pt,reqno]{amsart}

\usepackage[utf8]{inputenc}
\usepackage[T1]{fontenc}
\usepackage{amsmath,amssymb,amsthm,mathtools}
\usepackage{enumitem}
\usepackage[margin=1.15in]{geometry}
\usepackage[hidelinks]{hyperref}

\theoremstyle{plain}
\newtheorem{theorem}{Theorem}[section]
\newtheorem{proposition}[theorem]{Proposition}
\newtheorem{lemma}[theorem]{Lemma}
\newtheorem{corollary}[theorem]{Corollary}
\newtheorem{conjecture}[theorem]{Conjecture}
\theoremstyle{definition}
\newtheorem{definition}[theorem]{Definition}
\newtheorem{remark}[theorem]{Remark}

\newcommand{\T}{\mathbb{T}}
\newcommand{\R}{\mathbb{R}}
\newcommand{\N}{\mathbb{N}}
\newcommand{\Z}{\mathbb{Z}}
\newcommand{\E}{\mathbb{E}}
\newcommand{\HH}{\mathcal{H}}
\newcommand{\Had}{\mathbin{\circ}}     % Hadamard product
\newcommand{\Num}{\mathcal{N}}         % number operator
\newcommand{\Bm}{B^{H}}
\DeclareMathOperator{\spec}{spec}
\DeclareMathOperator{\tr}{tr}
\DeclareMathOperator{\diag}{diag}
\DeclareMathOperator{\Cov}{Cov}
\DeclareMathOperator{\Ran}{Ran}

\newcommand{\Gam}{\Gamma}

\title[An exact graded curvature tower for the fractional Laplacian]{An Exact
Graded Curvature Tower for the\\ Fractional Laplacian:\\ Operator Realization
and Wiener Chaos}
\author{Ramiro Fontes}
\address{Quijotic Research}
\email{ramirofontes@gmail.com}
\urladdr{https://www.quijotic.com}
\date{\today}

\subjclass[2020]{Primary 60J76, 47D07; Secondary 60H07, 39A70, 60G22}
\keywords{Bakry--\'Emery curvature; fractional Laplacian; carr\'e du champ;
iterated $\Gamma$-calculus; fractional Brownian motion; Wiener chaos; Hadamard
power; Cholesky factorization; number operator}

\begin{document}

\begin{abstract}
For the fractional Laplacian $-(-\Delta)^{\gamma/2}$ on the circle, $\gamma\in(0,2)$,
the carr\'e du champ of the associated jump generator coincides on positive
frequencies with the covariance of fractional Brownian motion with Hurst
parameter $H=\gamma/2$. We show that this is not a coincidence of the first
order but the seed of an exact, graded structure. The $k$-th iterated carr\'e
du champ $\Gam_k$ is realized as the squared norm of a \emph{diagonal tensor
lift} $e_n\mapsto e_n^{\otimes k}$ in the $k$-fold tensor power of the
fBM energy space, whose Gram matrix is the $k$-th Hadamard power $R_H^{\Had k}$
of the covariance. The $k$-th order Bakry--\'Emery curvature is therefore the
worst-case contraction ratio of the pencil $(R_H^{\Had k},R_H^{\Had(k-1)})$.

At the Brownian point $\gamma=1$ ($H=\tfrac12$, $R=[\min(n,m)]$) the whole tower
is solved in closed form. A single telescoping identity gives the Cholesky
factorization $R^{\Had k}=L\,D_k\,L^{\top}$ with $L$ the summation operator and
$D_k=\diag(j^{k}-(j-1)^{k})$ for \emph{every} $k$, from which
$R^{\Had(k-1)\,-1}R^{\Had k}$ is upper triangular with spectrum
$\{\,j^{k}-(j-1)^{k}\,\}_{j=1}^{N}$ --- the discrete derivative of $j\mapsto j^{k}$.
This recovers the odd integers at $k=2$, gives the uniform curvature lower bound
$\Gam_k\ge\Gam_{k-1}$ at every level (minimal eigenvalue $1$), and is generated
by a single number operator $\Num$ through $D_k=(\Num+1)^{k}-\Num^{k}$, so that
the universal constant $\kappa^{(k)}_1=1$ is the curvature of the Fock vacuum.
The determinant of the $k$-th curvature matrix \emph{decreases} monotonically
from $(2N-1)!!$ at $k=2$ to $N!$ as $k\to\infty$. The secular polynomial
$z^{k}-(z-1)^{k}$ has all of its zeros on the line $\Re z=\tfrac12$, so the
innovation variances are positive for a structural reason: a solvable instance
of ``critical line $\Rightarrow$ positivity,'' and the curvature spectral zeta at
$k=2$ is $(1-2^{-s})\zeta(s)$. The grading is the Wiener chaos of the Brownian
dual. Off the Brownian point the curvature deficit is transient: $\kappa^{(k)}
(\gamma,N)\to1$ exponentially in the chaos depth $k$, at the rate set by the
maximal off-diagonal fBM correlation, so every stable generator is
asymptotically Gaussian in its high-energy behavior.
\end{abstract}

\maketitle

%==============================================================================
\section{Introduction}\label{sec:intro}
%==============================================================================

\subsection*{The setting}
Let $\T=\R/2\pi\Z$ with normalized measure $\mu=dx/2\pi$, and let
$L_\gamma=-(-\Delta)^{\gamma/2}$, $\gamma\in(0,2)$, be the fractional Laplacian,
acting on Fourier modes by $L_\gamma e^{inx}=-|n|^{\gamma}e^{inx}$. It is the
generator of a symmetric pure-jump Markov semigroup with carr\'e du champ
\begin{equation}\label{eq:carre}
  \Gam(f,f)(x)=\tfrac12\int_{\R\setminus\{0\}}\!\bigl[f(x{+}z)-f(x)\bigr]^2
  \,\nu_\gamma(dz),\qquad \nu_\gamma(dz)=c_\gamma|z|^{-1-\gamma}\,dz,
\end{equation}
and iterated carr\'e du champ
$\Gam_2(f,f)=\tfrac12\bigl(L_\gamma\Gam(f,f)-2\Gam(f,L_\gamma f)\bigr)$, with the
higher levels defined recursively by
\begin{equation}\label{eq:gamma-rec}
  \Gam_{k}(f,f)=\tfrac12\bigl(L_\gamma\Gam_{k-1}(f,f)-2\Gam_{k-1}(f,L_\gamma f)\bigr),
  \qquad k\ge2 .
\end{equation}
The Bakry--\'Emery curvature condition $\mathrm{CD}(\kappa,\infty)$ at level $k$
asks for the largest $\kappa$ with $\Gam_k\ge\kappa\,\Gam_{k-1}$ pointwise; we
refer to~\cite{BGL} for the diffusion theory it generalizes.

Throughout we restrict to the space $\mathcal T_N^{+}$ of positive-frequency
trigonometric polynomials of degree at most $N$. On $\mathcal T_N^{+}$ the
Fourier kernel of $\Gam$ is, by a direct computation recorded in the companion
paper~\cite{FontesPartI},
\begin{equation}\label{eq:psi}
  \Psi_\gamma(\xi,\eta)=\tfrac12\bigl(|\xi|^{\gamma}+|\eta|^{\gamma}-|\xi-\eta|^{\gamma}\bigr)
  \;=\;R_H(\xi,\eta),\qquad H=\tfrac\gamma2,
\end{equation}
where $R_H(s,t)=\tfrac12\bigl(|s|^{2H}+|t|^{2H}-|s-t|^{2H}\bigr)$ is the
covariance of fractional Brownian motion (fBM) with Hurst parameter $H$. Thus the
\emph{base} curvature problem is a problem about fBM covariance matrices. This
paper is about what happens to the \emph{whole hierarchy}
\eqref{eq:gamma-rec}, and in particular about its exact solution at the
distinguished point $\gamma=1$.

\begin{remark}[Independence from the companion paper]
This paper may be read independently. We use \eqref{eq:psi} and the
operator-derivative framework as recalled in Section~\ref{sec:prelim}; the
connection with the fractional Bakry--\'Emery theory of~\cite{FontesPartI}
(arXiv:2602.19192) is motivational, and all statements proved below are
self-contained.
\end{remark}

\subsection*{Results}
Write $R_H=[R_H(n,m)]_{n,m=1}^{N}$, let $A^{\Had k}$ denote the $k$-th entrywise
(Hadamard) power of a matrix $A$, and let
$\kappa^{(k)}(\gamma,N)=\lambda_{\min}\!\bigl(R_H^{\Had k},R_H^{\Had(k-1)}\bigr)$
be the smallest generalized eigenvalue of the indicated pencil. Our results are
the following.

\begin{itemize}[leftmargin=1.4em]
\item \textbf{Realization} (Theorems~\ref{thm:opreal}--\ref{thm:tensor}). The
carr\'e du champ is the energy-space norm of the operator derivative
$D_{\Bm}=\delta_{\Bm}^{*}$ of a stochastic representative, and the $k$-th iterated
carr\'e du champ is the squared norm of the diagonal tensor lift
$e_n\mapsto e_n^{\otimes k}$, whose Gram matrix is $R_H^{\Had k}$. Hence the
graded curvature is a sequence of Hadamard-power pencils.

\item \textbf{The exact tower at $\gamma=1$} (Theorems~\ref{thm:cholesky}--%
\ref{thm:eig}, Proposition~\ref{prop:fock}). At the Brownian point
$R=[\min(n,m)]$, a single telescoping identity yields, for every $k$, the
Cholesky factorization $R^{\Had k}=L\,D_k\,L^{\top}$ with $L_{nj}=\mathbf 1_{n\ge j}$
and $D_k=\diag(j^{k}-(j-1)^{k})$, so that
\[
  \spec\!\bigl(R^{\Had(k-1)\,-1}R^{\Had k}\bigr)
  =\Bigl\{\,j^{k}-(j-1)^{k}\ :\ j=1,\dots,N\,\Bigr\},
\]
the discrete derivative of $j\mapsto j^{k}$. In particular $\kappa^{(2)}_j=2j-1$,
the level-$k$ curvature lower bound is $\Gam_k\ge\Gam_{k-1}$ uniformly
($\kappa^{(k)}_1=1$), and the tower is generated by a number operator
$\Num$ via $D_k=(\Num+1)^{k}-\Num^{k}$.

\item \textbf{Volume and a critical-line positivity model}
(Theorem~\ref{thm:det}, Theorem~\ref{thm:critline}). The determinant of the
$k$-th curvature matrix interpolates monotonically from $(2N-1)!!$ to $N!$. The
secular polynomial $z^{k}-(z-1)^{k}$ has all zeros on $\Re z=\tfrac12$, which
forces positivity of the innovation variances; the level-two spectral zeta is
$(1-2^{-s})\zeta(s)$.

\item \textbf{Off the Brownian point} (Theorem~\ref{thm:reg}). The curvature
deficit is transient in the chaos depth: $\kappa^{(k)}(\gamma,N)\to1$
exponentially, at the rate $\rho^{*}=\max_{n\neq m}|\rho_{nm}|<1$ given by the
maximal fBM correlation. Consequently every stable semigroup recovers Gaussian
energy decay at high chaos depth (Corollary~\ref{cor:gauss}).
\end{itemize}

The conceptual content is that the algebraic miracle behind the level-two phase
transition of~\cite{FontesPartI} is not special to level two. It is the
statement that the increments of Brownian motion are orthogonal, transported into
every tensor power of the energy space; the closed-form spectrum is the
innovation-variance sequence of the resulting Fock structure, and the Brownian
point is exactly the point at which those increments decouple.

%==============================================================================
\section{Preliminaries}\label{sec:prelim}
%==============================================================================

\subsection{The fBM energy space}
Sample fBM $\Bm$ at the integer times $\{1,\dots,N\}$. The associated
\emph{energy space} is $\HH=\R^{N}$ equipped with the inner product
\begin{equation}\label{eq:energy}
  \langle\mathbf u,\mathbf v\rangle_{\HH}=\sum_{n,m}u_n v_m\,R_H(n,m),
  \qquad\text{so that}\quad
  \langle e_n,e_m\rangle_{\HH}=\E[\Bm_n\Bm_m]=R_H(n,m).
\end{equation}
The stochastic integral (divergence) $\delta_{\Bm}(\mathbf u)=\sum_n u_n\Bm_n$ is
a bounded linear map $\HH\to L^2(\Omega)$, and the \emph{operator derivative} is
its Hilbert-space adjoint $D_{\Bm}:=\delta_{\Bm}^{*}$, characterized by
$\E[F\,\delta_{\Bm}(\mathbf u)]=\langle D_{\Bm}F,\mathbf u\rangle_{\HH}$. On the
first chaos one has $D_{\Bm}\bigl(\delta_{\Bm}(\mathbf b)\bigr)=\mathbf b$.

\subsection{Tensor powers}
The $k$-fold tensor power $\HH^{\otimes k}$ carries the inner product
\begin{equation}\label{eq:tensorip}
  \bigl\langle e_{n_1}\otimes\cdots\otimes e_{n_k},\,
  e_{m_1}\otimes\cdots\otimes e_{m_k}\bigr\rangle
  =\prod_{\ell=1}^{k}R_H(n_\ell,m_\ell).
\end{equation}
We single out the diagonal vectors $e_n^{\otimes k}$, for which
\eqref{eq:tensorip} gives the Gram matrix
\begin{equation}\label{eq:gram}
  \bigl\langle e_n^{\otimes k},e_m^{\otimes k}\bigr\rangle
  =R_H(n,m)^{k}=R_H^{\Had k}(n,m).
\end{equation}

\subsection{Stochastic representative}
For $f=\sum_{n=1}^{N}a_n e^{inx}\in\mathcal T_N^{+}$, set the phase-stripped
coefficient vector $\mathbf b(x)=(a_1e^{ix},\dots,a_Ne^{iNx})^{\top}$ and define
the stochastic representative
$\Phi_f(x):=\delta_{\Bm}(\mathbf b(x))=\sum_n b_n(x)\Bm_n$.

%==============================================================================
\section{Operator and tensor realization}\label{sec:real}
%==============================================================================

\begin{theorem}[Operator realization of the carr\'e du champ]\label{thm:opreal}
For $f\in\mathcal T_N^{+}$ and all $x\in\T$,
\begin{equation}\label{eq:opreal}
  \Gam(f,f)(x)=\bigl\|D_{\Bm}\Phi_f(x)\bigr\|_{\HH}^{2}.
\end{equation}
\end{theorem}

\begin{proof}
By the first-chaos identity $D_{\Bm}\Phi_f(x)=\mathbf b(x)$, so
$\|D_{\Bm}\Phi_f(x)\|_{\HH}^{2}=\langle\mathbf b(x),\mathbf b(x)\rangle_{\HH}
=\sum_{n,m}b_n(x)\overline{b_m(x)}\,R_H(n,m)$. Expanding
$\mathbf b(x)$ and using \eqref{eq:psi}, the right-hand side equals
$\sum_{n,m}a_n\overline{a_m}\,e^{i(n-m)x}\,\Psi_\gamma(n,m)$, which is precisely
the Fourier expansion of the jump form \eqref{eq:carre} on $\mathcal T_N^{+}$.
\end{proof}

\begin{theorem}[Tensor realization of the hierarchy]\label{thm:tensor}
Define the diagonal tensor lift $\iota_k:\R^N\to\HH^{\otimes k}$ by
$\iota_k(\mathbf v)=\sum_n v_n e_n^{\otimes k}$. Then for all $k\ge1$,
\begin{equation}\label{eq:tensorreal}
  \Gam_k(f,f)(x)=\bigl\|\iota_k(\mathbf b(x))\bigr\|_{\HH^{\otimes k}}^{2}
  =\sum_{n,m}b_n(x)\overline{b_m(x)}\,R_H^{\Had k}(n,m),
\end{equation}
and consequently
$\kappa^{(k)}(\gamma,N)=\lambda_{\min}\!\bigl(R_H^{\Had k},R_H^{\Had(k-1)}\bigr)$.
\end{theorem}

\begin{proof}
Each step of the recursion \eqref{eq:gamma-rec} multiplies the Fourier kernel by
one further factor of $\Psi_\gamma=R_H$ (the computation of~\cite{FontesPartI},
recalled in \eqref{eq:psi}), so the kernel of $\Gam_k$ on $\mathcal T_N^{+}$ is
the Hadamard power $R_H^{\Had k}$. By \eqref{eq:gram} this is exactly the Gram
matrix of the diagonal $k$-tensors, giving \eqref{eq:tensorreal}. The curvature
$\Gam_k\ge\kappa\,\Gam_{k-1}$ becomes
$\mathbf b^{*}R_H^{\Had k}\mathbf b\ge\kappa\,\mathbf b^{*}R_H^{\Had(k-1)}\mathbf b$
for all $\mathbf b$, i.e.\ $\kappa\le\lambda_{\min}(R_H^{\Had k},R_H^{\Had(k-1)})$.
\end{proof}

\begin{remark}[The grading is Wiener chaos]\label{rem:chaos}
The diagonal tensor $e_n^{\otimes k}$ is the image of $\Bm_n$ under the $k$-th
chaos map: if $\{H_k\}$ are the Hermite polynomials and $\Bm_n$ is normalized,
then $H_k(\Bm_n)$ spans the $k$-th Wiener chaos and
$\E[H_k(\Bm_n)H_k(\Bm_m)]=k!\,R_H(n,m)^{k}$~\cite{Nualart}. Thus
$R_H^{\Had k}=\tfrac1{k!}
\Cov\!\bigl(k\text{-th chaos}\bigr)$, and the Hadamard power appears for the
structural reason that the diagonal lift is the chaos embedding.
\end{remark}

%==============================================================================
\section{The exact tower at the Brownian point}\label{sec:tower}
%==============================================================================

Throughout this section $\gamma=1$, $H=\tfrac12$, and
$R:=R_{1/2}=[\min(n,m)]_{n,m=1}^{N}$.

\begin{lemma}[Telescoping factorization]\label{lem:telescope}
For every $k\ge1$,
\begin{equation}\label{eq:telescope}
  \min(n,m)^{k}=\sum_{j=1}^{N}\mathbf 1_{n\ge j}\,\mathbf 1_{m\ge j}\,
  \bigl(j^{k}-(j-1)^{k}\bigr).
\end{equation}
\end{lemma}

\begin{proof}
Both indicators are $1$ exactly when $j\le\min(n,m)$, so the sum telescopes:
$\sum_{j=1}^{\min(n,m)}\bigl(j^{k}-(j-1)^{k}\bigr)=\min(n,m)^{k}$.
\end{proof}

\begin{theorem}[Universal Cholesky factor]\label{thm:cholesky}
Let $L\in\R^{N\times N}$ be the summation operator $L_{nj}=\mathbf 1_{n\ge j}$
and $D_k=\diag\bigl(j^{k}-(j-1)^{k}\bigr)_{j=1}^{N}$. Then for every $k\ge1$,
\begin{equation}\label{eq:cholesky}
  R^{\Had k}=L\,D_k\,L^{\top}.
\end{equation}
The lower factor $L$ is independent of $k$; only the diagonal of innovation
variances $D_k$ depends on the chaos level. (For general background on Cholesky
and Hadamard-power structure see~\cite{HornJohnson}.)
\end{theorem}

\begin{proof}
$(L D_k L^{\top})_{nm}=\sum_j L_{nj}(D_k)_{jj}L_{mj}
=\sum_j\mathbf 1_{n\ge j}\mathbf 1_{m\ge j}(j^{k}-(j-1)^{k})$, which is
$\min(n,m)^{k}=R^{\Had k}(n,m)$ by Lemma~\ref{lem:telescope}.
\end{proof}

The geometric content of \eqref{eq:cholesky} is that the Brownian increments
$\Delta_j=e_j-e_{j-1}$ are orthonormal in $\HH_{1/2}$, and that their tensor
analogues
\begin{equation}\label{eq:incr}
  \delta_j^{(k)}:=e_j^{\otimes k}-e_{j-1}^{\otimes k}
  \qquad(e_0^{\otimes k}:=0)
\end{equation}
form an \emph{orthogonal} system in $\HH_{1/2}^{\otimes k}$ with squared norm
$\|\delta_j^{(k)}\|^{2}=j^{k}-(j-1)^{k}$; the columns of $L$ are the
change of basis $e_n^{\otimes k}=\sum_{j\le n}\delta_j^{(k)}$.

\begin{theorem}[Closed-form spectrum of the curvature tower]\label{thm:eig}
For every $k\ge2$ the $k$-th curvature matrix
$M^{(k)}:=R^{\Had(k-1)\,-1}R^{\Had k}$ is similar to a diagonal matrix and is
upper triangular in the standard basis up to the conjugation by $L^{\top}$, with
spectrum
\begin{equation}\label{eq:eig}
  \kappa_j^{(k)}=\frac{\|\delta_j^{(k)}\|^{2}}{\|\delta_j^{(k-1)}\|^{2}}
  =\frac{j^{k}-(j-1)^{k}}{\,j^{k-1}-(j-1)^{k-1}\,},\qquad j=1,\dots,N .
\end{equation}
In particular:
\begin{enumerate}[label=\textup{(\roman*)},leftmargin=2.2em]
\item $\kappa_j^{(2)}=2j-1$, recovering the odd integers $\{1,3,\dots,2N-1\}$;
\item $\kappa_1^{(k)}=1$ for every $k\ge2$ \textup(uniform curvature lower bound
$\Gam_k\ge\Gam_{k-1}$\textup);
\item $\kappa_j^{(k)}\downarrow j$ as $k\to\infty$ \textup(monotonically from above\textup).
\end{enumerate}
\end{theorem}

\begin{proof}
By Theorem~\ref{thm:cholesky}, $R^{\Had k}=LD_kL^{\top}$ and
$R^{\Had(k-1)}=LD_{k-1}L^{\top}$ with the same $L$. As $L$ is unit lower
triangular (hence invertible) and $D_{k-1}$ is positive (its entries
$j^{k-1}-(j-1)^{k-1}>0$),
\[
  M^{(k)}=\bigl(LD_{k-1}L^{\top}\bigr)^{-1}LD_kL^{\top}
  =L^{-\top}D_{k-1}^{-1}D_kL^{\top}.
\]
Thus $M^{(k)}$ is conjugate (by $L^{\top}$) to the diagonal matrix
$D_{k-1}^{-1}D_k$, whose entries are \eqref{eq:eig}; being a product of upper
triangular factors, $M^{(k)}$ is itself upper triangular with these numbers on
the diagonal. Item (i) is
$j^{2}-(j-1)^{2}=2j-1$; (ii) is $1^{k}-0^{k}=1$ for all $k$; and (iii) follows
since $\frac{j^{k}-(j-1)^{k}}{j^{k-1}-(j-1)^{k-1}}\to j$ as the leading terms
dominate.
\end{proof}

\begin{proposition}[Number-operator generator and the Fock vacuum]\label{prop:fock}
Let $\Num$ be the number operator $\Num\,\delta_j^{(k)}=(j-1)\,\delta_j^{(k)}$.
Then the innovation-variance operator is
\begin{equation}\label{eq:numop}
  D_k=(\Num+1)^{k}-\Num^{k},
\end{equation}
and the level-$k$ curvature acts on the $j$-th increment tensor as the ratio
$\bigl[(\Num+1)^{k}-\Num^{k}\bigr]\big/\bigl[(\Num+1)^{k-1}-\Num^{k-1}\bigr]$.
The universal eigenvalue $\kappa_1^{(k)}=1$ is the value at the vacuum $\Num=0$:
the curvature of the Fock vacuum is $1$ at every level.
\end{proposition}

\begin{proof}
On $\delta_j^{(k)}$ one has $\Num=j-1$, so $(\Num+1)^{k}-\Num^{k}=j^{k}-(j-1)^{k}
=\|\delta_j^{(k)}\|^{2}=(D_k)_{jj}$, which is \eqref{eq:numop}. The ratio formula
and the vacuum value $\frac{1^{k}-0^{k}}{1^{k-1}-0^{k-1}}=1$ are immediate.
\end{proof}

%==============================================================================
\section{Volume interpolation}\label{sec:vol}
%==============================================================================

\begin{theorem}[Determinant interpolation]\label{thm:det}
For every $k\ge2$,
\begin{equation}\label{eq:det}
  \det M^{(k)}=\prod_{j=1}^{N}\kappa_j^{(k)}
  =\prod_{j=1}^{N}\frac{j^{k}-(j-1)^{k}}{j^{k-1}-(j-1)^{k-1}} .
\end{equation}
At $k=2$ this is $\prod_{j=1}^{N}(2j-1)=(2N-1)!!$, and as $k\to\infty$ it
\emph{decreases} monotonically to $\prod_{j=1}^{N}j=N!$ \textup(note
$(2N-1)!!\ge N!$ for $N\ge2$\textup).
\end{theorem}

\begin{proof}
The determinant of $M^{(k)}$ is the product of its eigenvalues \eqref{eq:eig};
\eqref{eq:det} follows, and the $k=2$ value is the double factorial. Writing
$r=(j-1)/j\in[0,1)$ and $a_k=1-r^{k}$,
\[
  \kappa_j^{(k)}=j\,\frac{1-r^{k}}{1-r^{k-1}}=j\,\frac{a_k}{a_{k-1}} .
\]
The sequence $\{a_k\}$ is log-concave, since
$a_k^{2}-a_{k-1}a_{k+1}=r^{k-1}(1-r)^{2}\ge0$; hence the ratio $a_k/a_{k-1}$ is
nonincreasing in $k$, so each $\kappa_j^{(k)}$ decreases monotonically to its
limit $j$ (Theorem~\ref{thm:eig}(iii)). The product therefore decreases from
$(2N-1)!!$ to $\prod_j j=N!$.
\end{proof}

\begin{remark}
The two endpoints are the two natural ``volumes'': $(2N-1)!!=\E[X^{2N}]^{1/2}$-type
double factorial of the level-two (odd-integer) spectrum, and $N!$ the volume of
the fully decoupled Fock tower. The curvature ellipsoid shrinks from the former
to the latter as the chaos depth increases.
\end{remark}

%==============================================================================
\section{A critical-line positivity model}\label{sec:critline}
%==============================================================================

The innovation variances $D_k(j)=j^{k}-(j-1)^{k}$ are the integer samples of the
polynomial $p_k(z)=z^{k}-(z-1)^{k}$. Their positivity --- which is what makes the
whole tower a genuine curvature hierarchy --- is forced by a statement about the
\emph{location of the zeros} of $p_k$.

\begin{theorem}[All secular zeros lie on the critical line]\label{thm:critline}
For every $k\ge1$ the polynomial $p_k(z)=z^{k}-(z-1)^{k}$ has degree $k-1$, and
every one of its $k-1$ zeros satisfies $\Re z=\tfrac12$. Consequently $p_k$ is
real-rooted in the variable $\Re z$, $p_k(j)>0$ for every integer $j\ge1$, and
the sequence $\{D_k(j)\}_j$ is positive.
\end{theorem}

\begin{proof}
If $p_k(z)=0$ then $z^{k}=(z-1)^{k}$, hence $|z|=|z-1|$, which holds iff
$\Re z=\tfrac12$. The leading terms $z^{k}$ cancel, leaving degree $k-1$ and thus
exactly $k-1$ zeros, all on the line $\Re z=\tfrac12$. In particular every zero
has real part $\tfrac12<1$, so the real polynomial $x\mapsto p_k(x)$ has no zero
on $[1,\infty)$ and is of one sign there; since $p_k(1)=1>0$, we get
$D_k(j)=p_k(j)>0$ for all integers $j\ge1$.
\end{proof}

\begin{remark}[A solvable instance of the mechanism]
Theorem~\ref{thm:critline} is a fully solvable model of the implication
\emph{``all zeros on the critical line $\Rightarrow$ positivity of a moment/innovation
sequence''} --- the same shape of implication that a positivity proof of the
Riemann $\Xi$-function would require, here available in closed form because the
secular polynomial is elementary. We claim nothing about $\zeta$ beyond the
identity below; the point is structural, that the curvature tower is a setting in
which critical-line zeros and total nonnegativity are equivalent and both
provable.
\end{remark}

\begin{proposition}[Curvature spectral zeta]\label{prop:zeta}
The level-two curvature spectrum is the odd integers, and its spectral zeta
function is, for $\Re s>1$,
\begin{equation}\label{eq:zeta}
  \sum_{j\ge1}\bigl(\kappa_j^{(2)}\bigr)^{-s}
  =\sum_{j\ge1}(2j-1)^{-s}=\bigl(1-2^{-s}\bigr)\zeta(s).
\end{equation}
\end{proposition}

\begin{proof}
$\sum_{j\ge1}(2j-1)^{-s}=\sum_{n\ge1}n^{-s}-\sum_{j\ge1}(2j)^{-s}
=\zeta(s)-2^{-s}\zeta(s)$.
\end{proof}

%==============================================================================
\section{Off the Brownian point: chaos regularization}\label{sec:reg}
%==============================================================================

Away from $\gamma=1$ the increments are no longer orthogonal and the closed form
of Section~\ref{sec:tower} breaks. Nonetheless the curvature is recovered in the
limit of large chaos depth.

\begin{theorem}[Chaos regularization]\label{thm:reg}
Let $\rho_{nm}=R_H(n,m)/(n^{H}m^{H})$ be the fBM correlation and
$\rho^{*}=\max_{n\neq m}|\rho_{nm}|<1$. Then for all $\gamma\in(0,2)$ and all
$N\ge1$,
\begin{equation}\label{eq:reg}
  1-C_N(\rho^{*})^{\,k-1}\;\le\;\kappa^{(k)}(\gamma,N)\;\le\;1,
  \qquad k\ge2,\qquad C_N=2(N-1)N^{2H}.
\end{equation}
In particular $\kappa^{(k)}(\gamma,N)\to1$ exponentially in $k$ at rate
$\rho^{*}$.
\end{theorem}

\begin{proof}[Proof sketch]
Write $R_H^{\Had k}=D^{1/2}\bigl(\rho^{\Had k}\bigr)D^{1/2}$ with
$D=\diag(n^{2H})$ and $\rho=[\rho_{nm}]$ the correlation matrix ($\rho_{nn}=1$).
The off-diagonal entries of $\rho^{\Had k}$ are bounded by $(\rho^{*})^{k}$, so
$\rho^{\Had k}=I+E_k$ with $\|E_k\|\le(N-1)(\rho^{*})^{k}$. The generalized
eigenvalues of $(R_H^{\Had k},R_H^{\Had(k-1)})$ are those of
$(\rho^{\Had k},\rho^{\Had(k-1)})$ after the common conjugation by $D^{1/2}$ and
the diagonal scaling $\diag(n^{2H})$ between consecutive levels; a Weyl/Bauer--Fike
perturbation estimate against the identity gives the lower bound
\eqref{eq:reg} with the stated $C_N$. The upper bound $\kappa^{(k)}\le1$ is the
$\mathbf b=e_1$ test vector, for which $R_H^{\Had k}(1,1)=R_H^{\Had(k-1)}(1,1)=1$.
\end{proof}

\begin{corollary}[Asymptotic Gaussian energy decay]\label{cor:gauss}
Let $\mathcal E_k(f)=\int_\T\Gam_k(f,f)\,d\mu$ be the $k$-th energy. Along the
semigroup $P_t=e^{tL_\gamma}$ one has the dissipation identity
$\frac{d}{dt}\mathcal E_k(P_tf)=-2\,\mathcal E_{k+1}(P_tf)$ and hence the decay
rate $e^{-2\kappa^{(k+1)}(\gamma,N)\,t}$. By Theorem~\ref{thm:reg} this rate
approaches the Gaussian rate $e^{-2t}$ for every $\gamma\in(0,2)$ as $k\to\infty$:
every stable semigroup is asymptotically Gaussian in its high-energy behavior.
\end{corollary}

\begin{remark}[Optimality of the Brownian point]\label{rem:opt}
At level two, $\gamma=1$ is the unique maximizer of $\kappa^{(2)}(\gamma,N)$ over
$\gamma\in(0,2)$ (the phase transition of~\cite{FontesPartI}: cross-sign terms
vanish exactly at $\gamma=1$). Numerically the same holds at every level $k$ ---
$\gamma=1$ maximizes $\kappa^{(k)}(\gamma,N)$ --- but a level-by-level proof for
$k\ge3$ is not included here, and we record the higher-level optimality as a
conjecture supported by computation.
\end{remark}

%==============================================================================
\section{Outlook}\label{sec:outlook}
%==============================================================================

The structure isolated here is the positive-energy theory, valid on the fBM
range $H=\gamma/2\in(0,1)$. Two directions are natural. First, the Brownian point
$H=\tfrac12$ is the unique point at which the increment tensors decouple
\emph{exactly}; the closed-form tower, the number-operator generator, and the
critical-line positivity model are all manifestations of that decoupling, and it
would be of interest to characterize abstractly the Dirichlet forms whose
iterated carr\'e du champ admits such a universal Cholesky factor. Second, the
Hadamard-power realization $\Gam_k\leftrightarrow R_H^{\Had k}$ continues to make
formal sense past the endpoint $H=1$, where Schoenberg positivity of $R_H$ fails
and the carr\'e-du-champ form becomes indefinite; the resulting signed
Bakry--\'Emery geometry, governed by higher-order heat pseudo-processes, is the
subject of a companion paper. The present paper deliberately remains on the
positive side of that wall, where every statement is an equality or a genuine
curvature bound.

%==============================================================================
\begin{thebibliography}{9}

\bibitem{FontesPartI}
R.~Fontes,
\emph{Bakry--\'Emery curvature of the fractional Laplacian via fractional
Brownian covariance},
arXiv:2602.19192 [math.PR], 2026.

\bibitem{BGL}
D.~Bakry, I.~Gentil, and M.~Ledoux,
\emph{Analysis and Geometry of Markov Diffusion Operators},
Grundlehren der mathematischen Wissenschaften \textbf{348}, Springer, 2014.

\bibitem{Nualart}
D.~Nualart,
\emph{The Malliavin Calculus and Related Topics}, 2nd ed.,
Springer, 2006.

\bibitem{HornJohnson}
R.~A.~Horn and C.~R.~Johnson,
\emph{Matrix Analysis}, 2nd ed.,
Cambridge University Press, 2013.

\end{thebibliography}

\end{document}
